%% kapa: exam=1 round=23 year=2012 period=1 form=A subject=civil_law subject_ko=민법 seq=26 number=26
%% source: Q-Net 감정평가사 기출문제 (cst00309) artlSeq=5052578 file=raw/1008/5052578/2012년 제23회 감정평가사 제1차시험 기출문제(1교시 A형).pdf
%% conversion: kordoc 4.24.51 -> tools/gen1.py (auto); review=proofread-sonnet; flags=none
\begin{question}{26}
선의취득에 관한 설명으로 옳지 않은 것은? (다툼이 있으면 판례에 의함)
\begin{choices}
  \item 등록된 자동차는 선의취득의 대상이 아니다.
  \item 양수인의 선의ㆍ무과실의 기준시점은 물권행위가 완성된 때이다.
  \item 점유개정에 의한 인도에 대하여는 선의취득이 인정되지 않는다.
  \item 점유보조자가 본인의 의사에 반하여 본인의 물건을 처분하여 상대방이 선의취득한 경우에는 본인은 2년 내에 물건의 반환을 청구할 수 있다.
  \item 선의취득이 유효하게 성립한 경우에는 취득자는 임의로 선의취득의 효과를 거부하고 종전 소유자에게 동산을 반환받아 갈 것을 요구할 수 없다.
\end{choices}
\end{question}
\answer{4}
